
This study examined whether training corpus deduplication produces a uniform benchmark accuracy change or a contamination-proportional benchmark accuracy signature in Pythia models.

The primary finding is that the per-benchmark accuracy differential (dedup-Pile minus Pile) correlates positively with estimated 13-gram contamination rates across 16 observations (4 benchmarks x 4 model sizes, 160M to 6.9B parameters): Pearson r = 0.632, p = 0.0086; Spearman rho = 0.618, p = 0.0107; bootstrap 95\% CI = [0.297, 0.858]. MMLU, the benchmark with the highest evidence of contamination-driven inflation, shows Bonferroni-corrected significant accuracy reduction in dedup-Pile (t = -5.574, p = 0.0114).

A methodological contribution is the demonstration that token-count matching --- not step-matching --- is the correct confound-control approach for Pile/dedup-Pile comparisons. Step-matching introduces a volume confound that reduces measured contamination signal by Delta-r = 0.093 and introduces a non-contamination-related accuracy bias.

A negative result is also reported: the min-k\% memorization metric at Pythia-1B shows the opposite direction from prediction, indicating that the near-memorization pathway hypothesized to explain contamination-driven inflation does not manifest detectably via min-k\% at this model scale.

These findings suggest that MMLU accuracy reductions observed after corpus deduplication should be interpreted as contamination correction rather than model quality degradation, and that the benchmark-level direction and magnitude of deduplication's effects are predictable from corpus contamination estimates.

Future work should: (1) test min-k\% at Pythia-6.9B to determine whether the memorization signal manifests at larger scale; (2) re-run H-M3 with freshly computed H-M1 contamination estimates; (3) extend to 8--12 benchmarks to increase statistical power; (4) apply the framework to OLMo/Dolma for cross-family validation; and (5) compare against dedicated contamination mitigation baselines.

